\documentclass[12pt, a4paper]{exam}
\usepackage{fontspec}
\usepackage{xunicode}
\usepackage{xltxtra}
\usepackage{polyglossia}
\usepackage{graphicx}
\usepackage{tikz}
\usetikzlibrary{positioning, calc}
\usepackage{geometry}
\usepackage{enumitem}
\usepackage{multicol}
\usepackage{array}
\usepackage{lastpage}
\usepackage{amssymb}
\usepackage{multirow}
\usepackage{amsmath}

\setmainlanguage{thai}
\XeTeXlinebreaklocale "th"
\XeTeXlinebreakskip = 0pt plus 1pt

% Page Margins
\geometry{top=0.8in, bottom=1in, left=0.8in, right=0.8in}

% Fonts
\newfontfamily\thaifont[
  Scale=1.33,
  Script=Thai,
  Path = ../fonts/,
  Extension = .ttf,
  BoldFont = *-Bold,
  ItalicFont = *-Italic,
  BoldItalicFont = *-BoldItalic
]{THSarabunNew}

\setmainfont[
  Scale=1.33,
  Path = ../fonts/,
  Extension = .ttf,
  BoldFont = *-Bold,
  ItalicFont = *-Italic,
  BoldItalicFont = *-BoldItalic
]{THSarabunNew}

% Header and Footer
\pagestyle{headandfoot}
\extraheadheight{0.2in}

\newcommand{\examheader}{%
    ชื่อ-นามสกุล \makebox[3.0in]{\hrulefill} รหัสนักศึกษา \makebox[2.0in]{\hrulefill}
}


\firstpageheader{\examheader}{}{}
\runningheader{\examheader}{}{}

\firstpagefooter{\textit{ข้อสอบรายวิชา MTH234 Linear Algebra}}{}{\thepage}
\runningfooter{\textit{ข้อสอบรายวิชา MTH234 Linear Algebra}}{}{\thepage}

\begin{document}

\setcounter{page}{2}

\noindent \textbf{\underline{คำสั่ง}} : ให้นักศึกษาทำข้อสอบทุกข้อ โดยทำลงไปในข้อสอบ
\vspace{0.5cm}

\begin{questions}

% Question 1
\question \textbf{(8 คะแนน)} จงตอบคำถามต่อไปนี้
\begin{parts}
    \part ให้ $A = \begin{bmatrix} 2 & -1 & 5 \\ 3 & 0 & -2 \\ 1 & 4 & -3 \end{bmatrix}$ เป็นเมทริกซ์มาตรฐาน (standard matrix) สำหรับการแปลงเชิงเส้น (linear transformation) $T : \mathbb{R}^3 \to \mathbb{R}^3$ จงหาค่าของ $T\left(\begin{bmatrix} 1 \\ -2 \\ 3 \end{bmatrix}\right)$
    \vspace{\stretch{1}}
    
    \part กำหนดให้เมทริกซ์ $B = \begin{bmatrix} 1 & 0 & 0 & 0 \\ 0 & 3 & 0 & 0 \\ 0 & 0 & -2 & 1 \\ 0 & 0 & 0 & -2 \end{bmatrix}$ สามารถเขียนในรูปแบบของจอร์แดน (Jordan form) เป็น $\begin{bmatrix} M_1 & 0 \\ 0 & M_2 \end{bmatrix}$ จงระบุเมทริกซ์ $M_1$ และ $M_2$ สำหรับเมทริกซ์นี้
    \vspace{\stretch{1}}
    
    \newpage
    \part จงหาเวกเตอร์เจาะจงทั่วไป ซึ่งมีลำดับขั้นเท่ากับ 2 (generalized eigenvector of rank 2) ที่สัมพันธ์กับค่าเจาะจง (eigenvalue) $\lambda = 3$ ของเมทริกซ์ $C = \begin{bmatrix} 4 & -1 \\ 1 & 2 \end{bmatrix}$
    \vspace{\stretch{1}}
    
    \part กำหนดเมทริกซ์ $D = \begin{bmatrix} 7 & -4 \\ 5 & -2 \end{bmatrix}$ จงหาสมการลักษณะเฉพาะ (characteristic equation) และค่าเจาะจง (eigenvalues) ทั้งหมดของเมทริกซ์ $D$
    \vspace{\stretch{1}}
\end{parts}

\newpage

% Question 2
\question \textbf{(12 คะแนน)} กำหนดให้ $T : \mathbb{R}^4 \to \mathbb{R}^3$ นิยามโดย
\[ T\left(\begin{bmatrix} x_1 \\ x_2 \\ x_3 \\ x_4 \end{bmatrix}\right) = \begin{bmatrix} x_1 - x_2 + x_3 + x_4 \\ x_1 + 2x_2 - x_3 + x_4 \\ 2x_1 + x_2 + 2x_4 \end{bmatrix} \]
\begin{parts}
    \part จงแสดงว่า $T$ เป็นการแปลงเชิงเส้น (Show that $T$ is a linear transformation.)
    \vspace{\stretch{1}}
    \part จงหาฐานหลัก (basis) ของ $Ker(T)$ และค่าของ $nullity(T)$
    \vspace{\stretch{1}}
    \newpage
    \part จงหาฐานหลัก (basis) ของ $Range(T)$ และค่าของ $rank(T)$
    \vspace{\stretch{1}}
    \part จงตรวจสอบว่า $T$ เป็นฟังก์ชันหนึ่งต่อหนึ่ง (one-to-one) หรือไม่ และเป็นฟังก์ชันทั่วถึง (onto) หรือไม่
    \vspace{\stretch{1}}
\end{parts}

\newpage

% Question 3
\question \textbf{(12 คะแนน)} จงพิจารณาว่า เมทริกซ์ $A$ ที่กำหนดให้ต่อไปนี้เป็นเมทริกซ์ที่แปลงให้เป็นเมทริกซ์ทแยงมุมได้หรือไม่ (Determine if $A$ is diagonalizable.)
\[ A = \begin{bmatrix} 3 & -1 & 0 \\ -2 & 4 & 0 \\ 1 & 2 & 1 \end{bmatrix} \]
\begin{parts}
    \part จงหาค่าเจาะจง (eigenvalues) ทั้งหมดของ $A$
    \vspace{\stretch{1}}
    \part จงหาฐานหลักของปริภูมิเจาะจง (basis for eigenspace) สำหรับแต่ละค่าเจาะจง
    \vspace{\stretch{1}}
    \newpage
    \part หาก $A$ เป็นเมทริกซ์ที่แปลงให้เป็นเมทริกซ์ทแยงมุมได้ จงทำการหาเมทริกซ์ $P$ และเมทริกซ์ทแยงมุม (diagonal matrix) $D$ ซึ่งทำให้ $A = PDP^{-1}$ (ถ้าแปลงไม่ได้ให้ระบุเหตุผล)
    \vspace{\stretch{1}}
    \part จงแสดงวิธีการใช้เมทริกซ์ $P$ และ $D$ นี้ในการคำนวณหาเมทริกซ์ $A^5$
    \vspace{\stretch{1}}
\end{parts}

\newpage

% Question 4
\question \textbf{(11 คะแนน)} กำหนดให้
\[ Q(x,y) = 7x^2 + 13y^2 + 8xy \]
อยู่ในรูปแบบกำลังสอง (in the quadratic form)
\begin{parts}
    \part จงหาเมทริกซ์สมมาตร $A$ ที่สอดคล้องกับรูปแบบกำลังสองนี้
    \vspace{5cm}
    \part จงหาเมทริกซ์เชิงตั้งฉาก $P$ (orthogonal matrix) ที่ทำให้ $P^T A P = D$ เป็นเมทริกซ์ทแยงมุม
    \vspace{\stretch{1}}
    \newpage
    \part จงทำการเปลี่ยนตัวแปร $\vec{x} = P \vec{y}$ เพื่อให้รูปแบบกำลังสองของตัวแปรใหม่ ไม่มีเทอมผลคูณไขว้ปรากฏอยู่ (contains no cross-product terms) และระบุว่าเป็นสมการของกราฟรูปใด
    \vspace{\stretch{1}}
\end{parts}

\newpage

% Question 5
\question \textbf{(7 คะแนน)} การหมุนภาพในระนาบ (Rotation in 2D) ในระบบคอมพิวเตอร์กราฟิก คือการหมุนวัตถุรอบจุดกำเนิด (Rotation) ซึ่งถือเป็น Linear Transformation พื้นฐานที่สำคัญ หากกำหนดให้ $T : \mathbb{R}^2 \to \mathbb{R}^2$ เป็นฟังก์ชันการแปลงเชิงเส้นที่ทำการหมุนเวกเตอร์ $\vec{x}$ ทวนเข็มนาฬิกาเป็นมุม $\theta$ ไปยังจุดใหม่ $\vec{x}' = T(\vec{x})$ ความสัมพันธ์นี้สามารถเขียนในรูปสมการเมทริกซ์ได้เป็น
\[ T(\vec{x}) = R_\theta \vec{x} \]
โดยที่ $R_\theta$ คือ Rotation Matrix:
\[ R_\theta = \begin{bmatrix} \operatorname{\text{\fontfamily{cmr}\selectfont cos}}\theta & -\operatorname{\text{\fontfamily{cmr}\selectfont sin}}\theta \\ \operatorname{\text{\fontfamily{cmr}\selectfont sin}}\theta & \operatorname{\text{\fontfamily{cmr}\selectfont cos}}\theta \end{bmatrix} \]
ซึ่งเมทริกซ์นี้สร้างขึ้นจากการพิจารณาภาพของเวกเตอร์ฐานหลัก $\vec{e}_1 = \begin{bmatrix} 1 \\ 0 \end{bmatrix}$ และ $\vec{e}_2 = \begin{bmatrix} 0 \\ 1 \end{bmatrix}$ ภายใต้การหมุน

\begin{parts}
    \part จงหาเวกเตอร์ผลลัพธ์ $\vec{x}'$ จากการหมุนเวกเตอร์ $\vec{x} = \begin{bmatrix} 2 \\ 4 \end{bmatrix}$ ทวนเข็มนาฬิกาเป็นมุม $\theta = \dfrac{\pi}{3}$
    \vspace{\stretch{1}}
    
    \part จงหามุม $\theta$ (ในหน่วยเรเดียน) ที่ทำให้การหมุนภาพเปลี่ยนพิกัดจาก $\vec{x} = \begin{bmatrix} 1 \\ 1 \end{bmatrix}$ ไปยัง $\vec{x}' = \begin{bmatrix} -1 \\ 1 \end{bmatrix}$
    \vspace{\stretch{1}}
    
    \part จงพิสูจน์ว่า หากกำหนดฟังก์ชันการหมุนด้วยมุม $\theta = 2n\pi$ เมื่อ $n$ เป็นจำนวนเต็มใดๆ เมทริกซ์การแปลง $R_\theta$ จะสมมูลกับ เมทริกซ์เอกลักษณ์ (Identity Matrix)
    \vspace{\stretch{1}}
\end{parts}

\end{questions}
\end{document}
